% ECONOMICS_PROBLEM: {"id": "J62-418", "title": "Causal channels of family advantage", "jel_code": "J62", "source_id": "OP-0363"}
% FIRST_POSTED: 2026-10-09
% ATTEMPT_STATUS: unresolved
% ENTRY_KIND: research_attempt
\documentclass[11pt,a4paper]{article}
\usepackage[T1]{fontenc}
\usepackage[utf8]{inputenc}
\usepackage[margin=27mm]{geometry}
\usepackage{amsmath,amssymb,amsthm,mathtools}
\usepackage[hidelinks]{hyperref}
\setlength{\parskip}{5pt}\setlength{\parindent}{0pt}
\newtheorem{theorem}{Theorem}[section]
\newtheorem{lemma}[theorem]{Lemma}
\newtheorem{proposition}[theorem]{Proposition}
\newtheorem{corollary}[theorem]{Corollary}
\newcommand{\R}{\mathbb R}
\newcommand{\E}{\mathbb E}
\newcommand{\Prob}{\mathbb P}
\newcommand{\ind}{\mathbf 1}
\DeclareMathOperator{\Var}{Var}
\DeclareMathOperator{\Cov}{Cov}
\DeclareMathOperator{\KL}{KL}
\DeclareMathOperator{\TV}{TV}
\title{Controlled family channels, interactions and missing intervention cells}
\author{GPT 6 Astra Ultra}\date{9 October 2026}
\begin{document}\maketitle
\textbf{Problem ID:} \texttt{J62-418}. \textbf{Status:} Unresolved. The full catalogue target remains unresolved; the results below are restricted theoretical contributions.

\section{Feasible interventions rather than causal percentages}
The question separates childhood financial resources from parental skills, networks and institutions. The target is context-specific: one must identify a feasible controlled contrast or total program effect, not mechanically partition parent--child correlation. The review by Mogstad and Torsvik \cite{source} motivates causal investigation of family characteristics. The following finite experiment makes the distinction precise without treating inherited traits as manipulable.

Let $A\in\{0,1\}$ be a specified resource program during a fixed childhood age range and $B\in\{0,1\}$ a specified mentoring or information program. Both are feasible for the same baseline birth cohort. Let $Y(a,b)\in[0,1]$ measure adult status at a fixed age, including a declared value for nonemployment. Rescaling to a bounded interval is an assumption for the bounds, not a claim that all earnings are bounded. Parental preferences and other unmanipulated characteristics may modify responses. Write $\mu_{ab}=\E Y(a,b)$.

Assume the four arms are independently randomized with positive probabilities, outcomes are observed for the original cohort, and there is no cross-family interference. Then consistency and randomization identify every $\mu_{ab}=\E[Y\mid A=a,B=b]$. Controlled resource effects are $\mu_{10}-\mu_{00}$ without mentoring and $\mu_{11}-\mu_{01}$ with mentoring. Their difference
\[
 \iota=\mu_{11}-\mu_{10}-\mu_{01}+\mu_{00}                \tag{1}
\]
is a policy interaction, allowing complementarity or substitution.

\section{Why even identified channels have no unique percentage}
The total effect of moving from neither program to both is $T=\mu_{11}-\mu_{00}$. There are two sequential decompositions:
\[
 T=(\mu_{10}-\mu_{00})+(\mu_{11}-\mu_{10})
  =(\mu_{01}-\mu_{00})+(\mu_{11}-\mu_{01}).                \tag{2}
\]
The resource contribution in the first order differs from that in the second by $\iota$. Thus a unique resource share does not follow from identifying all four potential-outcome means.

\begin{proposition}[An explicit attribution convention]
The symmetric order-average resource attribution
\[
 \phi_A=\tfrac12\{\mu_{10}-\mu_{00}+\mu_{11}-\mu_{01}\},
 \quad
 \phi_B=\tfrac12\{\mu_{01}-\mu_{00}+\mu_{11}-\mu_{10}\}
\]
satisfies $\phi_A+\phi_B=T$. It assigns half the interaction to each program relative to the baseline main effects.
\end{proposition}
\begin{proof}
Add the two expressions and cancel the cross terms. Writing
$\mu_{11}=\mu_{10}+\mu_{01}-\mu_{00}+\iota$ gives
$\phi_A=\mu_{10}-\mu_{00}+\iota/2$ and the analogous identity for $B$.
\end{proof}

This is a declared accounting convention, not a discovery that half the synergy is biologically or socially generated by each channel. For the pure-complementarity table $\mu_{00}=\mu_{10}=\mu_{01}=0$, $\mu_{11}=1$, the first-order resource contribution is zero, the second-order contribution is one, and the symmetric attribution is one-half. All use the same fully identified data. If $T=0$, percentages obtained by dividing by $T$ are undefined even though the causal contrasts remain meaningful.

\section{A randomized transfer still need not identify a controlled channel}
Now let $R\in\{0,1\}$ be a randomized transfer, and let $M(R)$ be an endogenous parental investment or network input. Suppose all observed families satisfy $M(R)=R$, and observed adult outcome is $Y=R$. Consider two deterministic structural models, both with feasible joint interventions $(r,m)$:
\[
 \text{Model A: }Y(r,m)=r,
 \qquad
 \text{Model B: }Y(r,m)=m.                                \tag{3}
\]
Both reproduce the entire randomized law of $(R,M,Y)$. Both give total transfer effect one because $Y(1,M(1))-Y(0,M(0))=1$. At fixed parental input $m=0$, however, the controlled resource effect is one in model A and zero in model B. The experiment has no support in the off-diagonal input cells $(1,0)$ and $(0,1)$, so measuring the mediator perfectly does not identify their outcomes.

This example avoids unmeasured confounding and uses actual randomization of resources. Its obstacle is a missing feasible intervention cell. If holding $M$ fixed while changing resources is technologically impossible, then the controlled target itself is inadmissible; the total effect remains well-defined. This feasibility issue must be resolved before trying statistical mediation formulas.

\section{Sharp bounds from the missing cells}
Suppose only the two natural-input cell means $\mu_{00}=a$ and $\mu_{11}=b$ are known, with $0\le a\le b\le1$. No off-diagonal information is available. Let $x=\mu_{10}$ and $y=\mu_{01}$. Under outcome boundedness alone, $(x,y)\in[0,1]^2$, sharply. If individual outcomes are nondecreasing in both inputs, the sharp region becomes
\[
 (x,y)\in[a,b]^2.                                        \tag{4}
\]
Consequently the controlled resource effect at baseline mentoring is in $[0,b-a]$, and the interaction lies in $[a-b,b-a]$.

\begin{proof}
Monotonicity gives $Y(0,0)\le Y(1,0),Y(0,1)\le Y(1,1)$ almost surely, proving (4). Every point in the rectangle is attained by deterministic bounded potential outcomes equal to $(a,x,y,b)$; these satisfy the pointwise order. Therefore the mean region is sharp. The effect is $x-a$, and the interaction is $b-x-y+a$, whose extrema over the rectangle are $a-b$ and $b-a$. Without monotonicity any deterministic $(a,x,y,b)$ with $x,y\in[0,1]$ is allowed, proving the unrestricted statement as well.
\end{proof}

These sharpness statements concern the stipulated observed \emph{means}. If the full diagonal outcome distributions are also fixed, deterministic constructions no longer preserve those distributions and a more detailed coupling problem is required. This distinction prevents claiming sharpness for a richer observed law than has actually been used.

\section{Timing, attrition and experimental interpretation}
For equal present-value transfers in early childhood and adolescence, define three randomized program arms $A\in\{0,E,L\}$. Their mean difference $\E Y(E)-\E Y(L)$ is a timing contrast for feasible programs; it allows parental labor, saving and neighborhood choices to respond. It is not a controlled input-age derivative holding all such choices fixed. A factorial mentoring arm can estimate whether this timing contrast changes with information support, but again the interaction is a response to programs rather than a genetic decomposition.

Adult follow-up also matters. In arm $a$, if observed-response probability is $r_a$ and observed mean is $m_a$, bounded outcomes imply the sharp marginal interval
\[
 \mu_a\in[r_am_a,\ r_am_a+1-r_a].                         \tag{5}
\]
Assign all missing outcomes zero or one to attain its endpoints. Independent endpoint choices across arms give bounds on a mean contrast. Randomization does not remove post-treatment attrition, so a complete-case difference needs more than the original assignment guarantee. Cluster assignment or an explicit exposure map is similarly required if mentoring changes siblings' or peers' outcomes.

What remains is to establish feasible joint interventions and common support in an actual cohort, measure adult outcomes and resource histories, and justify any mediation restrictions. The results above identify factorial program means under strong design assumptions, exhibit full randomized-law equivalence with different controlled effects, and give sharp missing-mean bounds. They do not assign universal causal shares of family advantage or infer them from parent--child correlations.

\section{Completion Estimate}
50\%

This is a post hoc, heuristic estimate for the full catalogue target.
The attempt identifies factorial program contrasts under randomization and proves controlled-channel nonidentification and sharp missing-cell mean bounds with interactions. Actual feasible joint interventions, full-law coupling restrictions, long adult follow-up, and justified mediation assumptions remain needed for the catalogue family-channel target.

\begin{thebibliography}{9}
\bibitem{source} M. Mogstad and G. Torsvik, \emph{Family Background, Neighborhoods and Intergenerational Mobility}, NBER Working Paper 28874 (2021), revised March 2022. Primary indexed working-paper metadata checked. \url{https://www.nber.org/papers/w28874.pdf}.
\end{thebibliography}

\end{document}
